\documentclass{article}

\begin{document}
    \begin{center}
        \Large \textbf{Rabbit Challenge 0: Das Einrichten der Entwicklungsumgebung ist Teil des Projekts } \\ [6pt]
        \normalsize
        Autor: Maximilian Jakob Maag B.Sc. \\
        Version: 1.0
    \end{center}

    \section{Intro}
    Die Challenge richtet sich an die Rabbit Fraktion. Aufgabe ist das Einrichten einer lokalen IDE.
    Sie müssen in der Lage sein VMs zu deployen.
    \\ \\
    Die Konfiguration der Elemente erfolgt dann über Ansible.

    \section{Aufgaben}
    Am Ende dieser Challenge sollten Sie ein funktionierendes Setup besitzen, um die folgenden Challenges zu meistern.

    \subsection{Aufgabe 1}
    Erzeugen Sie ein SSH Schlüsselpaar.

    \subsection{Aufgabe 2}
    Implementieren Sie ein Terraform Modul, dass Sie immer wieder verwenden können um schnell und einfach eine VM zu deployen.
    Ihre VM muss folgende Eigenscahften besitzen:
    \begin{itemize}
        \item Ubuntu 22.04
        \item 2 vCPUs
        \item 4 GB RAM
        \item 20 GB Festplatte
        \item Port 22 offen
        \item Ihr SSH Public Key im authorized_keys File
        \item SSH Key Ihres Huntsman Konkurenten im authorized_keys File
    \end{itemize}

    \subsection{Aufgabe 3}
    Erstellen SIe ein Ansible Playbook, dass die VM mit einem Webserver ausstattet und diesen startet.

    \şubsection{Aufgabe 4}
    Erstellen Sie ein weiteres Ansible Playbook. Das Playbook legt eine Datei mit dem Namen: rabbit-was-here.txt in /var/www/rabbit an.
    Notieren Sie in dieser Datei Ihren Namen um für dieses Deployment auf dem Leaderboard einen Punkt zu erhalten.
    
    \subsection{Aufgabe 5}
    Richten Sie Ihr Notebook so ein, dass Sie ein Ansible Playbook auf einer VM ausführen können.
    Sie dürfen davon ausgehen, dass auf der VM Ihr SSH Public Key im authorized_keys File liegt.

\end{document}